\chapter*{L'approche séquence vers séquence}
\label{chap:seq2seq}
On introduit l'étude, on aborde les difficultés du problème de base et qu'on étudie ici un problème similaire simplifié.
\section*{Description du problème}
On introduit le problème simplifié, en insistant sur les modifications avec le problème de base et les apports de cette étude.
\subsection*{Encodage du temps}
\label{seq:timeencoding}
\begin{equation}
\label{eq:timeencodingin}
a = b + 1
\text{Ici on montre comment le temps est encodé dans les séquences}
\end{equation}

\begin{equation}
\label{eq:timeencodingout}
a = b + 1
\text{Ici on montre comment le temps est encodé dans les séquences}
\end{equation}
\section*{Architectures employées}
On introduit les architectures étudiées pour cette partie, il faut justifier ces choix. Il s'agit d'un transformer mais on pourra ajouter des résultats sur les RNNs permettant de les disqualifier. Les CNNs ne s'adapte pas au traitement seq2seq.
\section*{Suite : (ce n'est pas un titre de section}
Dans la suite on présentera comme dans la soumission à CAID.
\section*{Traitement de séquence entière}
\section*{Fenêtrage}
\section*{Discussion}
\label{seq:discseq2seq}


\chapter{L'approche séquence vers séquence}
Le chapitre présente la principale méthode de résolution de la problématique : une approche de transduction séquence ver séquence. Cette approche est la plus naturelle du point de vue de l'environnement virtuel puisque le module DCI génère toutes les impulsions avant que les modules suivantes ne traitent la séquence d'interceptions entière. Ce chapitre présente les travaux que constituent un article soumis à la conférence CAID. Il est articulé en trois sections. Nous commençons en exposant à nouveau les motivations de cette approches, en formulant nos attendes et nos réserves par rapport aux chances de succès. Ensuite nous présentons des travaux concernant le traitement des séquences entières, approches dont les limites sont liées aux tailles des séquences traitées. Finalement, nous présentons des travaux sur un mécanisme de fenêtrage, permettant de traiter les séquences longues en plusieurs fois, et son intégration dans un processus de prédiction de séquences.

\section{Motivations}
Les Transformers excellent 
